\chapter{기준계, 빛, 그리고 시공간의 기하학}

\section{상대성이론으로 다시 시작하는 강의}

이번 학기에는 고전물리학으로 돌아가지만, 공간과 시간이 서로 독립적이고 보편적이라고 보는 뉴턴 역학의 관점으로 돌아가지는 않는다. 이번에 다룰 주제는 특수상대성이론과 고전장 이론, 특히 전자기학이다. 양자역학은 사용하지 않는다. 앞선 역학 강의에서는 작용 원리와 해밀턴 및 라그랑주 기술을 다루었고, 양자 강의에서는 기본 원리를 다루었다. 다만 조화 진동자조차 나중에 다시 다루기로 하고 남겨 두어야 했다. 빛도 대부분 다루지 않았다. 빛은 본질적으로 상대론적이므로, 이제 그 빠진 내용을 보충한다.

수강생 대부분은 특수상대성이론을 전에 접해 보았다. 그래도 손을 들어 조사해 보니 처음부터 다시 세우는 편이 낫다는 의견이 더 많다. 그것이 유익한 선택이다. 로런츠 변환을 외워야 할 공식으로 제시해서는 안 된다. 공간과 시간의 조작적 측정이 갈릴레이 운동학과 관측된 빛의 전파 법칙을 동시에 만족시킬 수 없기 때문에 로런츠 변환이 도출되어야 한다.

\subsection{물리적 장치로서의 좌표}

데카르트 좌표계를 생각해 보자. 원점과 방향이 정해진 세 축, 그리고 좌표 \(x,y,z\)가 있다. 이를 구체적으로 구현하는 한 방법은 공간 전체에 미터 자로 된 단단한 격자를 놓는 것이다. 단순히 장소를 지정하는 것이 아니라 사건을 특정하려면 격자의 각 지점에 시계도 있어야 한다. 그러면 사건은 다음과 같이 표시된다.
\begin{equation}
(x,y,z,t).
\end{equation}
시계들은 명시된 절차에 따라 동기화해야 한다. 이 요건이 곧 구성에서 까다로운 부분이 된다.

\begin{definition}
기준계란 공간 좌표 격자와 그 격자 전체에 분포한 동기화 시계들의 체계다. 관성계란 회전하거나 가속하지 않는 기준계로, 그 안에서는 힘을 받지 않는 모든 입자가 일정한 속도로 운동한다.
\end{definition}

서로 등속 운동하는 두 기준계는 같은 사건에 서로 다른 좌표를 부여한다. 움직이는 기준계에 고정된 점의 좌표 \(x'\)는 일정하지만, 그 점의 비프라임 좌표는 변한다. 뉴턴 물리학에서 \(x\) 방향의 운동에 대한 공간 좌표 관계는 다음과 같다.
\begin{equation}
x'=x-vt.
\label{eq:sr1-galilean-space}
\end{equation}
눈에 덜 띄는 뉴턴 역학의 가정은 잘 맞춰진 모든 시계가 어디서나 동기화되어 있다는 것이다.
\begin{equation}
t'=t.
\label{eq:sr1-absolute-time}
\end{equation}

\subsection{상대성 원리}

상대성이론은 아인슈타인과 함께 시작된 것이 아니다. 갈릴레이와 뉴턴도 물리 법칙이 모든 관성계에서 같은 형태를 갖는다는 원리를 이미 사용했다. 뉴턴의 법칙 \(F=ma\), 역제곱 중력 법칙, 쿨롱 힘은 실험실이 등속 운동한다고 해서 새로운 형태를 갖지 않는다.

친숙한 사고실험은 창문이 밀폐된 열차 칸이다. 열차가 매끄럽게 움직이고 창밖을 볼 수 없다면, 공을 저글링하는 동작은 정지한 방에서와 똑같다. 공은 평소처럼 올라갔다 내려와 손으로 돌아온다. 열차 칸 안에서만 하는 실험으로는 열차의 등속도를 알아낼 수 없다. 가속이나 회전은 내부에서 감지할 수 있지만, 일정한 병진 운동은 감지할 수 없다.

아인슈타인은 이 오래된 원리에 물리 법칙 하나를 더한다. 진공에서 빛은 다음 속도로 전파된다.
\begin{equation}
c \simeq 1.86\times 10^5\ {\rm miles/s}
   \simeq 3.00\times 10^8\ {\rm m/s}.
\end{equation}
이것이 자연 법칙이고 모든 관성계에서 법칙이 같다면, 모든 관성 관측자는 같은 \(c\)를 측정해야 한다.

\section{갈릴레이 변환이 실패하는 이유}

한 공간 차원만으로도 충돌을 드러내기에 충분하다. 가로축을 \(x\), 세로축을 \(t\)로 하는 시공간 도표를 그려 보자. 원점에서 방출되어 오른쪽으로 진행하는 빛 펄스는 다음 궤적을 따른다.
\begin{equation}
x=ct.
\label{eq:sr1-light-right}
\end{equation}
속력 \(v\)로 오른쪽으로 움직이는 기준계의 원점은 다음 궤적 위에 놓인다.
\begin{equation}
x=vt,
\qquad
x'=0.
\label{eq:sr1-moving-origin}
\end{equation}

\begin{figure}[htbp]\centering
\includegraphics[width=0.90\textwidth]{lecture_01_figure_01.png}\caption{오른쪽으로 진행하는 빛의 광선과 움직이는 원점 \(x=vt\). 프라임 관측자는 이 원점을 \(x'=0\)라고 표시한다.}\lectureframe{00:15:24}\end{figure}

빛의 광선에 갈릴레이 규칙 \eqref{eq:sr1-galilean-space}--\eqref{eq:sr1-absolute-time}을 적용해 보자.
\begin{align}
x'&=x-vt=(c-v)t,\\
t'&=t.
\end{align}
그러면 프라임 관측자는 다음과 같이 추론하게 된다.
\begin{equation}
x'=(c-v)t',
\end{equation}
따라서 오른쪽으로 진행하는 빛의 속력은 \(c-v\)가 된다. 왼쪽으로 진행하는 광선 \(x=-ct\)에 대해서는
\begin{equation}
x'=-(c+v)t',
\end{equation}
속력의 크기가 \(c+v\)이다.

\begin{figure}[htbp]\centering
\includegraphics[width=0.90\textwidth]{lecture_01_figure_02.png}\caption{갈릴레이 치환은 오른쪽으로 진행하는 빛의 광선에 대해 \(c-v\)를 예측하므로, 기준계에 무관한 빛의 속도와 모순된다.}\lectureframe{00:17:23}\end{figure}

일반적인 속도 덧셈은 느린 역학적 운동에서 훌륭한 근사이지만, 빛을 측정해도 예측된 \(c-v\)와 \(c+v\)가 관측되지는 않는다. 빛의 공리를 바꿀 것이 아니라 변환 법칙을 재고해야 한다. 핵심적으로 의심해야 할 것은 겉보기에는 무해한 명제 \(t'=t\)이다. 이 명제는 절대적 동시성의 개념을 가정한다.

\section{멀리 떨어진 시계 동기화하기}

서로 다른 위치의 두 시계가 같은 시각을 가리킨다고 말하려면 실험이 필요하다. 두 시계의 중간 지점에 세 번째 관측자를 두고, 미터 자로 두 시계까지의 거리가 같음을 확인한다. 양 끝의 시계가 정오를 가리킬 때 섬광을 내도록 한다. 섬광이 중간 지점에 동시에 도착하면 두 시계를 동기화된 것으로 정의한다.

이것은 물리 위에 덧씌우는 철학적 정의가 아니다. 시계, 거리, 빛 신호로 수행하는 재현 가능한 조작이다.

이제 한 관측자가 중간 지점을 지나가며 움직인다고 하자. 정지한 중간 지점에 동시에 도착한 양 끝의 섬광은 움직이는 관측자에게는 동시에 도착하지 않는다. 빛이 이동하는 동안 관측자는 한쪽 섬광 쪽으로 다가가고 다른 쪽 섬광에서는 멀어지기 때문이다. 따라서 정지한 시계들로 동시에 판정한 사건들은, 움직이는 기준계에서 똑같이 타당한 동기화 절차를 따르면 동시가 아니다.

비프라임 기준계에서 동시 사건들의 직선은 다음과 같다.
\begin{equation}
t=\text{constant}.
\end{equation}
이에 대응하는 프라임 직선은 \(x\)-\(t\) 도표에서 기울어진다. 변환 공식을 묻기 전에 빛 신호로 그 기울기를 구성할 수 있다.

\section{빛으로 그린 움직이는 기준계}

\(c=1\)로 두자. 시간을 초로 잴 때 거리를 광초로 잴 수도 있고, 정한 길이 단위를 빛이 가로지르는 데 걸리는 시간으로 시간을 잴 수도 있다. 어느 관례를 따르든 빛의 광선은 다음과 같다.
\begin{equation}
x=\pm t,
\label{eq:sr1-light-units}
\end{equation}
따라서 두 축의 눈금을 같은 비율로 잡으면 \(45^\circ\) 직선이 된다.

먼저 빛의 광선을 그린다. 그런 다음 움직이는 기준계에서 정지해 있는 관측자들의 무리를 그린다.
\begin{equation}
x=vt,\qquad x=vt+1,\qquad x=vt+2.
\label{eq:sr1-moving-family}
\end{equation}
관측자들은 자기 관성 격자에서 서로 간격을 일정하게 유지하므로, 이들의 세계선은 서로 평행하다.

\begin{figure}[htbp]\centering
\includegraphics[width=0.90\textwidth]{lecture_01_figure_03.png}\caption{\(45^\circ\)의 빛 광선과 움직이는 좌표 격자의 평행한 세계선 \(x=vt\), \(x=vt+1\), \(x=vt+2\).}\lectureframe{00:29:55}\end{figure}

\subsection{프레드, 메리, 시모어}

프레드는 \(x=vt\)에, 메리는 \(x=vt+1\)에, 시모어는 \(x=vt+2\)에 있다고 하자. 원점 \(O\)에서 프레드는 오른쪽으로 진행하는 빛 신호를 메리에게 보낸다. 시모어도 왼쪽으로 진행하는 신호를 보내 두 섬광이 같은 사건 \(a\)에서 메리에게 도달하게 한다. 움직이는 기준계의 중점 동기화 규칙에 따르면 시모어가 신호를 방출한 사건 \(b\)은 프레드가 \(O\)에서 신호를 방출한 사건과 동시다. 그러므로 \(O\)과 \(b\)을 지나는 직선은 \(t'=0\)이다.

사건 \(a\)은 두 식 모두를 만족한다.
\begin{equation}
x=t
\end{equation}
그리고
\begin{equation}
x=vt+1.
\end{equation}
따라서
\begin{equation}
t_a=x_a=\frac{1}{1-v}.
\label{eq:sr1-event-a}
\end{equation}

\(a\)을 지나는 왼쪽 방향의 빛 광선은 다음 관계를 만족한다.
\begin{equation}
x+t=x_a+t_a=\frac{2}{1-v}.
\label{eq:sr1-ab-light}
\end{equation}
사건 \(b\)도 시모어의 세계선 위에 있으므로,
\begin{equation}
x-vt=2.
\label{eq:sr1-seymour-line}
\end{equation}
\eqref{eq:sr1-ab-light}와 \eqref{eq:sr1-seymour-line}을 풀면 다음을 얻는다.
\begin{align}
t_b&=\frac{2v}{1-v^2},\\
x_b&=\frac{2}{1-v^2}.
\label{eq:sr1-event-b}
\end{align}
따라서 \(O\)에서 \(b\)까지 잇는 직선의 기울기는 다음과 같다.
\begin{equation}
\frac{t_b}{x_b}=v,
\end{equation}
그러므로 움직이는 기준계의 동시성 직선은 다음과 같다.
\begin{equation}
t=vx,
\qquad
t'=0.
\label{eq:sr1-moving-simultaneity}
\end{equation}

\begin{figure}[htbp]\centering
\includegraphics[width=0.93\textwidth]{lecture_01_figure_04.png}\caption{프레드--메리--시모어의 동기화 구성. 도착 시각이 같은 빛 신호를 이용하면 움직이는 기준계에서 \(O\)와 \(b\)가 동시임을 확인할 수 있으며, 그 결과 \(t=vx\)을 얻는다.}\lectureframe{00:44:38}\end{figure}

이제 프라임 기준계의 두 축은 각각 다음과 같은 평행선이다.
\begin{align}
x'=0 &\quad\Longleftrightarrow\quad x=vt,\\
t'=0 &\quad\Longleftrightarrow\quad t=vx.
\end{align}
두 축은 오른쪽으로 진행하는 빛의 직선 \(x=t\)을 기준으로 서로 거울 대칭이다. 이 기하학적 구성은 이미 보편적 동시성을 대신하고 있다.

\begin{classroomqa}
\audiencequestion{\(v=1\) 또는 \(v=-1\)일 때 이 구성에는 어떤 일이 일어나는가?}
\lecturerresponse{\eqref{eq:sr1-event-a}--\eqref{eq:sr1-event-b}의 분모가 0이 되어 해당 교점들은 무한대로 멀어진다. \(\lvert v\rvert=1\)일 때 질량을 가진 관측자는 빛의 광선을 따라 움직이게 되어 통상적인 중점 교환을 완료할 수 없다. 이러한 특이 거동은 \(c\)가 한계 속도임을 보여 주는 최초의 기하학적 경고이다. \lecturetimestamp{00:40:43}}
\end{classroomqa}

\section{로런츠 변환 유도하기}

\(x'=0\)일 때 \(x=vt\)이므로 \(x'\)에는 \(x-vt\) 인자가 포함되어야 한다. 마찬가지로 \(t'=0\)일 때 \(t=vx\)이므로 \(t-vx\) 인자가 포함되어야 한다. 시공간의 균질성으로 변환은 선형이므로 다음과 같이 쓴다.
\begin{align}
x'&=(x-vt)f(v),\\
t'&=(t-vx)g(v).
\label{eq:sr1-linear-ansatz}
\end{align}

\begin{figure}[htbp]\centering
\includegraphics[width=0.90\textwidth]{lecture_01_figure_05.png}\caption{선형 변환은 두 영점 조건으로 구성되며, 속도에 따라 달라지는 스케일 인자 \(f(v)\)와 \(g(v)\)를 포함한다.}\lectureframe{00:49:35}\end{figure}

오른쪽으로 진행하는 빛의 광선 \(x=t\)을 따라가면, 빛의 속도 불변성에 따라 \(x'=t'\)이어야 한다. 식~\eqref{eq:sr1-linear-ansatz}은 그러면 다음과 같이 된다.
\begin{align}
x'&=(1-v)x f(v),\\
t'&=(1-v)x g(v),
\end{align}
따라서
\begin{equation}
f(v)=g(v).
\label{eq:sr1-f-equals-g}
\end{equation}
그리하여 공통 스케일 인자 하나만 남는다.
\begin{align}
x'&=(x-vt)f(v),\\
t'&=(t-vx)f(v).
\label{eq:sr1-direct-unscaled}
\end{align}
이는 아인슈타인의 첫 상대성이론 논문에서 사용한 간결한 논증과 사실상 같다. 조작적으로 정의한 동시성과 불변인 빛의 속도로 변환을 결정한다.

\subsection{상호성이 스케일을 결정한다}

\emph{실제로 움직이는} 특별한 관성계는 없다. 프라임 기준계가 비프라임 기준계에 대해 \(v\)로 움직인다면, 비프라임 기준계는 프라임 기준계에 대해 \(-v\)로 움직인다. 따라서 역변환도 \(v\)의 부호만 반대로 바뀐 같은 형태여야 한다.
\begin{align}
x&=(x'+vt')f(v),\\
t&=(t'+vx')f(v).
\label{eq:sr1-inverse-unscaled}
\end{align}

\begin{figure}[htbp]\centering
\includegraphics[width=0.90\textwidth]{lecture_01_figure_06.png}\caption{상호성에 따라 상대 속도의 부호를 반대로 바꾸면 같은 스케일 인자를 유지한 역변환을 얻는다.}\lectureframe{00:55:12}\end{figure}

\eqref{eq:sr1-direct-unscaled}을 \eqref{eq:sr1-inverse-unscaled}의 첫 번째 식에 대입하자.
\begin{align}
x
&=\left[(x-vt)f(v)+v(t-vx)f(v)\right]f(v)\\
&=x(1-v^2)f(v)^2.
\end{align}
이 합성 변환은 항등 변환이어야 하므로,
\begin{equation}
(1-v^2)f(v)^2=1,
\end{equation}
따라서
\begin{equation}
f(v)=\pm\frac{1}{\sqrt{1-v^2}}.
\label{eq:sr1-factor-roots}
\end{equation}

\begin{figure}[htbp]\centering
\includegraphics[width=0.90\textwidth]{lecture_01_figure_07.png}\caption{순변환과 역변환을 합성하면 공통 인자 \(f(v)\)의 크기가 결정된다.}\lectureframe{00:59:31}\end{figure}

\begin{classroomqa}\audiencequestion{\eqref{eq:sr1-factor-roots}에서 음의 제곱근을 버리는 이유는 무엇인가?}\lecturerresponse{상대 속도가 0일 때 변환은 연속적으로 항등 변환에 가까워져야 한다. 즉, \(x'=x\)이고 \(t'=t\)이어야 한다. 음의 가지를 택하면 대신 \(x'=-x\)과 \(t'=-t\)이 되어 두 좌표의 방향이 모두 뒤집힌다. 이를 택하면 기준계가 연속적으로 움직이는 대신 갑작스러운 반사가 도입된다. \lecturetimestamp{00:59:27}}\end{classroomqa}

양의 가지를 택하고 \(c=1\) 단위를 사용하면 로런츠 변환은 다음과 같다.
\begin{align}
x'&=\frac{x-vt}{\sqrt{1-v^2}},&
t'&=\frac{t-vx}{\sqrt{1-v^2}}.
\label{eq:sr1-lorentz-c1}
\end{align}

\begin{figure}[htbp]\centering
\includegraphics[width=0.90\textwidth]{lecture_01_figure_08.png}\caption{빛의 속도를 1로 둔 단위계에서 완성된 로런츠 변환.}\lectureframe{01:02:05}\end{figure}

\subsection{일반 단위로 되돌리기}

차원 분석을 통해 일반 단위를 복원할 수 있다. 분모에는 무차원 비 \(v^2/c^2\)만 들어갈 수 있고, 시간 방정식의 \(vx\)는 \(vx/c^2\)가 되어야 한다.
\begin{align}
x'&=\gamma(x-vt),&
t'&=\gamma\left(t-\frac{vx}{c^2}\right),
\label{eq:sr1-lorentz-dimensional}\\
\gamma&=\frac{1}{\sqrt{1-v^2/c^2}}.
\label{eq:sr1-gamma}
\end{align}
역변환은 다음과 같다.
\begin{align}
x&=\gamma(x'+vt'),&
t&=\gamma\left(t'+\frac{vx'}{c^2}\right).
\label{eq:sr1-lorentz-inverse-dimensional}
\end{align}

\begin{figure}[htbp]\centering
\includegraphics[width=0.90\textwidth]{lecture_01_figure_09.png}\caption{\(c\)를 복원한 로런츠 공식. \(v^2/c^2\)은 스케일을 결정하고 \(vx/c^2\)는 위치를 시간에 섞는다.}\lectureframe{01:05:01}\end{figure}

\(v\ll c\)일 때,
\begin{equation}
\gamma=1+\frac{1}{2}\frac{v^2}{c^2}+O\!\left(\frac{v^4}{c^4}\right),
\end{equation}
변환은 다음과 같이 된다.
\begin{equation}
x'\approx x-vt,
\qquad
t'\approx t.
\end{equation}
지상에서의 속도에서는 이 비들이 대개 매우 작으므로 뉴턴 운동학이 잘 들어맞는다. \(c\)에 가까워지면 보정 효과가 커진다.

운동 방향은 \(x\)축을 따라가도록 정했다. 상호성을 어기지 않고서는 횡방향 측정 막대로 두 관성계를 구별할 수 없으므로,
\begin{equation}
y'=y,
\qquad
z'=z.
\end{equation}

\section{움직이는 막대의 길이 측정}

길이는 측정하는 기준계에서 동시에 측정한 두 끝점 사이의 거리다. 사진에 찍힌 물체의 겉보기 각크기를 뜻하지는 않는다. 빛의 이동 시간 지연으로 별도의 광학 문제가 생기기 때문이다.

비프라임 기준계에서 정지해 있는 막대의 양 끝점이 다음과 같다고 하자.
\begin{equation}
x=0,
\qquad
x=1.
\end{equation}
프라임 관측자는 프라임 기준계의 한 동시성 단면에서 두 끝점에 해당하는 사건을 측정한다. 예를 들어 \(t'=0\)를 택하자. 먼 쪽 끝점 \(x=1\)에서
\begin{equation}
t'=0
\quad\Longrightarrow\quad
t=\frac{vx}{c^2}.
\end{equation}
잠시 \(c=1\) 기준계로 돌아가면, 이 사건은 \(t=v\)이고
\begin{equation}
x'=\frac{x-vt}{\sqrt{1-v^2}}
   =\frac{1-v^2}{\sqrt{1-v^2}}
   =\sqrt{1-v^2}.
\end{equation}

\begin{figure}[htbp]\centering
\includegraphics[width=0.90\textwidth]{lecture_01_figure_10.png}\caption{비프라임 기준계에서 정지한 막대의 세계선 \(x=0\)과 \(x=1\). 움직이는 관측자는 하나의 \(t'\) 단면에서 두 세계선을 비교해야 한다.}\lectureframe{01:11:14}\end{figure}

\begin{figure}[htbp]\centering
\includegraphics[width=0.90\textwidth]{lecture_01_figure_11.png}\caption{프라임 기준계에서 동시에 측정한 막대 양 끝점의 거리는 \(\sqrt{1-v^2}\)만큼 수축한다. 여기서는 \(c=1\) 단위를 쓴다.}\lectureframe{01:16:03}\end{figure}

일반 단위와 임의의 정지 길이 \(L_0\)로 되돌리면,
\begin{equation}
L=\frac{L_0}{\gamma}
 =L_0\sqrt{1-\frac{v^2}{c^2}}.
\label{eq:sr1-length-contraction}
\end{equation}

\begin{classroomqa}\audiencequestion{\(x=1\)에 막대의 먼 쪽 끝이 있다면 그 좌표도 \(x'=1\)이어야 하지 않나요?}\lecturerresponse{아니다. 막대는 비프라임 기준계에서 정지해 있으므로, 그 기준계에서 \(x=0\)와 \(x=1\)는 막대의 고정된 양 끝점을 나타낸다. 먼 쪽 끝점의 프라임 좌표는 시간에 따라 변한다. 측정 단면 \(t'=0\)에서의 그 좌표를 로런츠 변환으로 계산한다. \lecturetimestamp{01:14:24}}\end{classroomqa}

\begin{classroomqa}\audiencequestion{각 관성 관측자가 상대방의 막대가 더 짧다고 일관되게 말할 수 있는 이유는 무엇인가?}\lecturerresponse{두 관측자는 서로 다른 끝점 사건 쌍을 측정한다. 비프라임 관측자는 \(t\)가 같은 사건을 요구하고, 프라임 관측자는 \(t'\)가 같은 사건을 요구한다. 동시성이 기준계에 따라 달라지므로, 이 두 끝점 사건 쌍은 같은 사건 쌍이 아니다. 따라서 상호성은 모순을 낳지 않는다. \lecturetimestamp{01:16:28}}\end{classroomqa}

조작적 의미를 분명히 해 두자. 측정 기준계 곳곳에 동기화된 시계와 보조자를 둔다. 미리 정한 현지 시각에 각 보조자가 가까이 지나가는 막대 끝점의 위치를 기록한다. 동시에 기록한 두 위치를 빼면 그 기준계에서 막대의 길이가 나온다. 막대가 어떻게 보이는지는 전혀 관계없다.

\begin{classroomqa}\audiencequestion{\(c=1\)일 때 단위 막대의 길이는 반드시 광초여야 하나요?}\lecturerresponse{시간을 초로 측정할 때에만 그렇다. 길이는 미터로 쓰고 시간을 1미터를 빛이 가로지르는 데 걸리는 시간으로 잴 수도 있다. \(c=1\)로 두면 시간과 길이의 단위가 서로 연결되지만, 어떤 물리적 단위를 먼저 선택해야 하는지를 정하는 것은 아니다. \lecturetimestamp{01:21:17}}\end{classroomqa}

\section{움직이는 시계와 쌍둥이 문제}

이제 프라임 기준계에서 정지한 시계를 따라가 보자. 그 시계의 세계선을 따라
\begin{equation}
x'=0.
\end{equation}
그 시계가 \(t'=1\)를 가리키는 사건을 원점 사건과 비교하자. 역변환을 적용하면
\begin{equation}
t=\gamma\left(t'+\frac{vx'}{c^2}\right)=\gamma.
\end{equation}
따라서 움직이는 시계에서 경과한 시간 1단위는 비프라임 좌표 시간 \(\gamma\) 단위에 해당한다.
\begin{equation}
\Delta t=\gamma\,\Delta\tau,
\qquad
\Delta\tau=\frac{\Delta t}{\gamma}.
\label{eq:sr1-time-dilation}
\end{equation}

\begin{figure}[htbp]\centering
\includegraphics[width=0.90\textwidth]{lecture_01_figure_12.png}\caption{\(x'=0\)를 따라가는 시계는 같은 두 사건 사이에서 정지한 좌표계의 시각보다 적은 시간을 기록한다. 즉, \(t=\gamma t'\)이다.}\lectureframe{01:26:42}\end{figure}

이 효과는 시계라고 부르는 기계 장치에만 국한되지 않고 모든 물리적 시계에 적용된다. 손목시계, 원자 전이, 신진대사, 생물학적 노화 모두 같은 고유 시간을 따른다. 뉴욕 거리에서 산 싸구려 시계도 롤렉스보다 성능이 떨어진다고 해서 상대성이론을 피할 수는 없다. 정확한 시각을 유지하지 못할 수는 있지만, 어떤 과정을 이용하든 그 과정에는 같은 운동학이 적용된다.

\begin{classroomqa}\audiencequestion{각 관성 관측자는 상대방의 시계가 느리게 간다고 보는데, 여행한 쌍둥이가 더 젊은 모습으로 돌아오는 이유는 무엇인가?}\lecturerresponse{왕복 여행은 전역적으로 동등한 두 관성 세계선을 비교하는 상황이 아니다. 여행자는 방향을 바꾸어 다른 관성계에 들어가야 하지만, 집에 남은 쌍둥이는 하나의 관성 세계선에 머문다. 방향 전환으로 교환 대칭성이 깨지므로 누적된 고유 시간은 서로 다를 수 있다. \lecturetimestamp{01:26:37}}\end{classroomqa}

\begin{classroomqa}\audiencequestion{방향을 바꾸는 동안의 가속 자체가 나이 차이를 일으키나요?}\lecturerresponse{가속은 대칭성이 깨지는 지점을 표시하지만, 나이 차이는 전체 경로를 따라 고유 시간을 적분해 계산한다. 방향 전환을 짧게 하거나 관성 운동 구간 사이의 변화로 다루어도, 서로 다른 시공간 경로에서는 여전히 고유 시간이 다르다. \lecturetimestamp{01:29:20}}\end{classroomqa}

\section{불변 기하학}

시간 지연 도표를 유클리드 기하학으로 읽으면 모순처럼 보인다. 기울어진 경로가 빗변처럼 보이지만 수직 경로보다 고유 시간이 짧다. 이를 해결하는 것은 대수적 속임수가 아니다. 시공간의 계량이 다르기 때문이다.

유클리드 평면에서 회전은 \(x\)와 \(y\)를 바꾸지만 다음 양은 보존한다.
\begin{equation}
x'^2+y'^2=x^2+y^2.
\label{eq:sr1-euclidean-invariant}
\end{equation}
로런츠 변환이 \(t^2+x^2\)을 보존하는지 확인해 볼 수 있다. 보존하지 않는다. \(c=1\) 단위에서는
\begin{align}
t'^2
&=\frac{t^2-2vtx+v^2x^2}{1-v^2},\\
x'^2
&=\frac{x^2-2vtx+v^2t^2}{1-v^2}.
\end{align}
두 식을 더하면 계수가 맞지 않게 된다. 대신 빼면 혼합항들이 상쇄되어 다음을 얻는다.
\begin{align}
t'^2-x'^2
&=
\frac{t^2-2vtx+v^2x^2-x^2+2vtx-v^2t^2}{1-v^2}\\
&=t^2-x^2.
\label{eq:sr1-interval-invariance}
\end{align}

\begin{figure}[htbp]\centering
\includegraphics[width=0.90\textwidth]{lecture_01_figure_13.png}\caption{변환된 좌표의 제곱을 서로 빼면 교차항이 상쇄되어 불변량 \(t^2-x^2\)이 드러난다.}\lectureframe{01:38:07}\end{figure}

좌표는 기준계에 따라 달라지는 성분이다. 하지만 다음 조합은
\begin{equation}
\Delta s^2=c^2\Delta t^2-\Delta x^2
\label{eq:sr1-spacetime-interval}
\end{equation}
모든 관성계에서 같다. 공간 차원이 더 많으면
\begin{equation}
\Delta s^2
=c^2\Delta t^2-\Delta x^2-\Delta y^2-\Delta z^2.
\end{equation}

\begin{classroomqa}\audiencequestion{빗변이 더 길다는 통상적인 기하학 법칙을 버려도 되는 이유는 무엇인가요?}\lecturerresponse{그 법칙은 양의 정부호인 유클리드 계량에 속한다. 여기서 좌표는 자, 시계, 빛 신호를 이용해 조작적으로 정의했으며, 그 좌표의 변환은 제곱의 차를 보존한다. 민코프스키 기하학은 역설을 피하려고 가정한 것이 아니라, 이러한 측정에서 발견한 것이다. \lecturetimestamp{01:43:21}}\end{classroomqa}

\subsection{고유 시간}

시계의 세계선 위에 있는 두 사건을 생각하고 그 시계의 정지 기준계를 택하자. 그 기준계에서는 \(\Delta x'=0\)이므로 불변성에 따라
\begin{equation}
c^2\Delta\tau^2
=c^2\Delta t^2-\Delta x^2,
\end{equation}
또는
\begin{equation}
\Delta\tau
=\sqrt{\Delta t^2-\frac{\Delta x^2}{c^2}}.
\label{eq:sr1-proper-time}
\end{equation}
이것이 두 사건 사이에서 시계가 실제로 기록하는 시간이다.

\begin{figure}[htbp]
\centering
\includegraphics[width=0.90\textwidth]{lecture_01_figure_14.png}
\caption{시계의 정지계에서는 공간적 분리가 사라지므로, 불변 간격은 고유 시간의 제곱이 됩니다.}
\lectureframe{01:41:28}
\end{figure}

\begin{classroomqa}\audiencequestion{고유 시간은 \(t^2-x^2/c^2\)인가요, 아니면 그 제곱근인가요?}\lecturerresponse{제곱의 차가 고유 시간의 제곱이다. 공통 원점으로부터 잰 분리에 대해 시계가 읽는 경과 시간은 \(\tau=\sqrt{t^2-x^2/c^2}\)이다. 처음 설명에서 빠진 제곱근을 복원해야 한다. \lecturetimestamp{01:51:02}}\end{classroomqa}

\subsection{시간꼴, 광꼴, 공간꼴 분리}

서로 다른 두 사건의 간격이 0일 수 있다. \(c=1\) 단위에서는
\begin{equation}
\Delta t^2-\Delta x^2=0
\quad\Longleftrightarrow\quad
\Delta x=\pm\Delta t.
\end{equation}
이러한 사건들은 빛의 광선으로 연결된다.

\begin{figure}[htbp]
\centering
\includegraphics[width=0.90\textwidth]{lecture_01_figure_15.png}
\caption{빛꼴로 분리된 사건들은 \(x=\pm t\) 위에 놓이며, 시공간 간격과 빛과 같은 경로를 따른 고유 시간은 영이 됩니다.}
\lectureframe{01:46:40}
\end{figure}

빛이 보통의 시계를 운반할 수 있다면, 방출과 흡수 사이에 고유 시간 눈금은 하나도 흐르지 않는다. 이것은 빛의 정지 기준계를 뜻하지 않는다. 모든 광꼴 구간의 고유 시간이 0이라는 극한적 진술이다.

빛의 광선은 시공간을 세 가지 불변 분류로 나눈다.
\begin{align}
c^2\Delta t^2-\Delta x^2&>0
&&\text{timelike},\\
c^2\Delta t^2-\Delta x^2&=0
&&\text{null},\\
c^2\Delta t^2-\Delta x^2&<0
&&\text{spacelike}.
\end{align}
물리적 시계는 시간꼴로 분리된 사건들을 연결할 수 있으며, 양의 간격으로 시계의 고유 시간을 정한다. 아광속 물체는 공간꼴로 분리된 사건들을 연결할 수 없다. 공간꼴 사건 쌍에 대해서는 두 사건이 동시에 일어나는 기준계를 택해 고유 거리
\begin{equation}
\ell_{\rm proper}
=\sqrt{\Delta x^2-c^2\Delta t^2}.
\end{equation}
를 정의할 수 있다.

\begin{figure}[htbp]\centering
\includegraphics[width=0.90\textwidth]{lecture_01_figure_16.png}\caption{빛원뿔은 시계로 연결할 수 있는 시간꼴 사건 쌍과 아광속 물체로 연결할 수 없는 공간꼴 사건 쌍을 가른다.}\lectureframe{01:48:47}\end{figure}

\begin{classroomqa}\audiencequestion{불변량을 고유 시간이라고 부르고 고유 길이라고 부르는 때는 각각 언제인가요?}\lecturerresponse{시간꼴 사건 쌍에는 두 사건이 같은 장소에서 일어나는 정지 기준계가 있으며, 이때 간격은 고유 시간을 준다. 공간꼴 사건 쌍에는 두 사건이 같은 시각에 일어나는 기준계가 있으며, 부호가 반대인 불변량은 두 사건 사이의 고유 거리를 준다. 광꼴 분리는 이 두 경우의 경계다. \lecturetimestamp{01:47:22}}\end{classroomqa}

\begin{classroomqa}\audiencequestion{\(c\)는 \(c=1\)로 정하지 않은 단위계에서 어디에 들어가나요?}\lecturerresponse{불변량은 \(c^2\Delta t^2-\Delta x^2\)이며, 차원은 길이의 제곱이다. 이를 \(c^2\)로 나누면 시간의 제곱 차원을 갖는 \(\Delta\tau^2=\Delta t^2-\Delta x^2/c^2\)를 얻는다. \lecturetimestamp{01:49:56}}\end{classroomqa}

\section{로런츠 대칭성과 그 역사적 의미}

\begin{classroomqa}\audiencequestion{로런츠 변환을 회전으로 볼 수 있나요?}\lecturerresponse{그렇다. 다만 보통의 유클리드 회전은 아니다. 로런츠 변환은 \(t^2-x^2\)을 보존하는 쌍곡 회전이다. 역사적으로는 허수각만큼의 회전이라고 설명할 수 있고, 현대에는 급속도를 이용해 실수로 매개화한다. \lecturetimestamp{01:51:40}}\end{classroomqa}

아인슈타인이 택한 접근은 전자기학과 밀접하게 연결되어 있었다. 훗날 회고에 따르면, 그는 십 대 때 이미 빛의 파동을 쫓아가면 어떻게 될지 궁금해했다. 갈릴레이식 속도 뺄셈에 따르면 파동이 정지한 것으로 보이겠지만, 맥스웰 방정식에는 진공에서 정지해 진동하는 빛의 파동에 해당하는 해가 없다. 이 충돌은 측정 막대의 특이한 역학적 성질이 아니라 운동학을 가리켰다.

\begin{classroomqa}\audiencequestion{아인슈타인의 논증은 주로 마이컬슨--몰리 같은 실험에 대한 대응이었나요?}\lecturerresponse{실험적 배경은 역사적으로 중요했지만, 아인슈타인 자신의 회고에서는 맥스웰 방정식의 구조와 빛줄기를 쫓아가는 사고실험이 강조되었다. 그가 특정 간섭계 실험 결과를 얼마나 알고 있었거나 논증에 얼마나 의존했는지는 역사학자들 사이에서 논쟁 중이다. 따라서 명확한 물리적 요점은 맥스웰 이론과 갈릴레이 변환 사이의 대칭성 충돌이다. \lecturetimestamp{01:52:04}}\end{classroomqa}

맥스웰 방정식은 갈릴레이 변환이 아니라 로런츠 변환에 대해 불변이다. 로런츠와 피츠제럴드도 에테르의 존재를 유지하고 수축을 에테르 속 운동에 따른 역학적 변형으로 해석하면서 관련된 수축 공식을 이미 찾아냈다. 그 관점에서는 이 효과가 추가 보정을 수반하는 근사적인 물질의 반응일 수 있었다.

아인슈타인의 결정적인 진전은 상대성 원리와 빛의 속도 불변성을 정확한 운동학적 원리로 삼은 것이었다. 그러면 길이 수축과 시간 지연은 적절한 모든 막대와 시계에 적용된다. 이것들은 에테르가 가하는 특수한 힘이 아니라 관성계 사이의 기하학적 관계를 나타내기 때문이다. 다음 과제는 이 기하학 안에서 입자의 운동을 정식화하는 것이다.
